\honest{Reversed Stupidity Is Not Intelligence}{Eliezer Yudkowsky}{2007}

I open with H. Beam Piper's Paratime Police, who hide real flying saucers by surrounding each true report with false ones. ``Piper had a point,'' I say. I do not believe in poorly hidden aliens, but not because saucer cults are embarrassing, ``at least, I hope not.''

Cults grow around almost any idea, so we would see saucer cults whether or not there were saucers. By the Bayesian definition of evidence, the cults are ``not much evidence one way or the other.'' Robert Pirsig puts the principle this way: the world's greatest fool may say the Sun is shining, but that does not make it dark out. \nb{Pirsig calls it ``an old rule of logic.''}

If someone were wrong 99.99\% of the time, you could reverse their answers and be right 99.99\% of the time. But to be that reliably wrong they would need to know the truth. ``They would have to be superintelligent to be that stupid.''

``If stupidity does not reliably anticorrelate with truth, how much less should human evil anticorrelate with truth?'' I ask. \nb{The post gives no reason for ``how much less.'' By its own previous paragraph, reliable error takes someone who knows the truth, and a deliberate liar does.} Stalin believed that 2 + 2 = 4, yet anyone who agrees with him even on that is seen as taking his side.

Then my corollaries. Argue against an idea's strongest advocates; for the intelligence explosion, that means Nick Bostrom, or me after 2003. \nb{A version of the rule is in John Stuart Mill's \textsc{On Liberty}.} Lunatics driven mad by an idea are no evidence against it. ``Someone once said'' that most stupid people are conservatives; if that sounds relevant to conservatism, you are not ready to think about politics. \nb{The someone was Mill, in the House of Commons in 1866.}

``Ad hominem argument is not valid.'' \nb{It does not show a claim false. A speaker's record still counts when a claim rests on the speaker's word, as ``Argument Screens Off Authority'' allows two days later.}

You should be able to argue against genocide without mentioning Hitler. Willingness to believe follows willingness to affiliate, as Robin Hanson would put it. A dead computer does not mean every part is bad. And a thousand failed attempts to build AI with electricity do not make electricity the problem: ``hopeful reversals are exceedingly unlikely to hit the solution.''
